\documentclass[11pt,a4paper]{article}
\usepackage[T1]{fontenc}
\usepackage[utf8]{inputenc}
\usepackage{amsmath,amssymb,amsthm}
\usepackage[margin=1in]{geometry}
\usepackage[colorlinks=true,linkcolor=blue,urlcolor=blue]{hyperref}
\theoremstyle{plain}
\newtheorem{theorem}{Theorem}
\newtheorem{lemma}[theorem]{Lemma}
\newtheorem{proposition}[theorem]{Proposition}
\newtheorem{corollary}[theorem]{Corollary}
\theoremstyle{definition}
\newtheorem{remark}[theorem]{Remark}

\title{The empirical closed form for OEIS A240335, proved}
\author{Adrian Perez Fontelles\\ \small Independent researcher}
\date{1 September 2026}
\begin{document}
\maketitle

\begin{abstract}
OEIS A240335 carries an empirical closed form for $a(n)$ --- not a recurrence relating
consecutive terms but an explicit formula. It is true, and it is decidable rather than
empirical. The entry counts arrays subject to a condition on a bounded window of consecutive lines, so the lines of an array are the
vertices of a finite digraph, an array is a walk in it, and $a(n)$ is a walk count on
$S=1992$ vertices. Any function of the form $\sum_j p_j(n)\lambda_j^{\,n}$ satisfies the
monic linear recurrence whose characteristic polynomial is
$\prod_j(x-\lambda_j)^{\deg p_j+1}$; here that polynomial is $q(x)=(x-1)^{31}$, of degree
$31$. Two things are then checked exactly: that the walk count satisfies the same
recurrence from some index on, which the Cayley--Hamilton theorem reduces to a finite
computation, and that the two sequences agree at $31$ consecutive indices past that
point. A monic recurrence determines every later term from its predecessors, so agreement
there is agreement for good.
\end{abstract}

\noindent\small 2020 Mathematics Subject Classification. 05A15, 11B37, 15A18.\normalsize

\section{The conjecture}

OEIS A240335 is ``Number of nX4 0..3 arrays with no element equal to two plus the sum of elements to its left or two plus the sum of elements above it or one plus the sum of the elements diagonally to its northwest, modulo 4.''. It has offset $1$ and begins
\[
11, 125, 1543, 14456, 110327, 716770, 4106515, 21225132,\ \dots
\]
The entry states, as a conjecture contributed by its author and never marked settled:
\begin{quote}
Empirical polynomial of degree 30 (see link above)

Empirical: a(n) = (1/4144575934565485291192320000000)*n\^{}30 + (1/36840674973915424810598400000)*n\^{}29 + (113/13450956379928552669184000000)*n\^{}28 + (31/473429106539928354816000000)*n\^{}27 + (1419401/18148115750697253601280000000)*n\^{}26 - (247991/93067260259985915904000000)*n\^{}25 + (62270911/106604316297802049126400000)*n\^{}24 - (4757068873/130294164363980282265600000)*n\^{}23 + (814130487823/254923365059961421824000000)*n\^{}22 - (6095076727/31530410025969254400000)*n\^{}21 + (25984521815341/2427841571999632588800000)*n\^{}20 - (368487574420439/809280523999877529600000)*n\^{}19 + (2270389462764494647/174421249777868341248000000)*n\^{}18 + (42913688146108121/430669752537946521600000)*n\^{}17 - (97938993233169055441/2052014703269039308800000)*n\^{}16 + (414425876629714148377/97714985869954252800000)*n\^{}15 - (200356654541928443577659/755070345358737408000000)*n\^{}14 + (572120899038507158656891/42593711789467238400000)*n\^{}13 - (2587261421801842724447494619/4498647618705201561600000)*n\^{}12 + (8194719386341948927635005737/386247522818123366400000)*n\^{}11 - (324590832832412561336280617460011/477981309487427665920000000)*n\^{}10 + (8586770985635419192193990436961/455220294749931110400000)*n\^{}9 - (1316018047637568195385804280059579/2908351883124559872000000)*n\^{}8 + (4502755798196345951073817056013433/484725313854093312000000)*n\^{}7 - (5918556804941401097645230792739739487/36758336300602076160000000)*n\^{}6 + (61595135975260835139015152068553953/26636475580146432000000)*n\^{}5 - (14923705744317372817320729322890599/556944489403061760000)*n\^{}4 + (3513221073658897687484396660166227/14586641389127808000)*n\^{}3 - (147007268067241018958241836628521/93256746094512000)*n\^{}2 + (299407299371860260636969629/44790183900)*n - 13783108115112102 for n>35
\end{quote}
As of the ``Last modified'' line on the live entry (Jun 02 2025 09:49:35, revision 7) this is still
recorded as empirical, and nothing on the entry records it as proved. Write
\[
f(n)\;=\;\frac{n^{30}}{4144575934565485291192320000000} + \frac{n^{29}}{36840674973915424810598400000} + \frac{113 n^{28}}{13450956379928552669184000000} + \frac{31 n^{27}}{473429106539928354816000000} + \frac{1419401 n^{26}}{18148115750697253601280000000} - \frac{247991 n^{25}}{93067260259985915904000000} + \frac{62270911 n^{24}}{106604316297802049126400000} - \frac{4757068873 n^{23}}{130294164363980282265600000} + \frac{814130487823 n^{22}}{254923365059961421824000000} - \frac{6095076727 n^{21}}{31530410025969254400000} + \frac{25984521815341 n^{20}}{2427841571999632588800000} - \frac{368487574420439 n^{19}}{809280523999877529600000} + \frac{2270389462764494647 n^{18}}{174421249777868341248000000} + \frac{42913688146108121 n^{17}}{430669752537946521600000} - \frac{97938993233169055441 n^{16}}{2052014703269039308800000} + \frac{414425876629714148377 n^{15}}{97714985869954252800000} - \frac{200356654541928443577659 n^{14}}{755070345358737408000000} + \frac{572120899038507158656891 n^{13}}{42593711789467238400000} - \frac{2587261421801842724447494619 n^{12}}{4498647618705201561600000} + \frac{8194719386341948927635005737 n^{11}}{386247522818123366400000} - \frac{324590832832412561336280617460011 n^{10}}{477981309487427665920000000} + \frac{8586770985635419192193990436961 n^{9}}{455220294749931110400000} - \frac{1316018047637568195385804280059579 n^{8}}{2908351883124559872000000} + \frac{4502755798196345951073817056013433 n^{7}}{484725313854093312000000} - \frac{5918556804941401097645230792739739487 n^{6}}{36758336300602076160000000} + \frac{61595135975260835139015152068553953 n^{5}}{26636475580146432000000} - \frac{14923705744317372817320729322890599 n^{4}}{556944489403061760000} + \frac{3513221073658897687484396660166227 n^{3}}{14586641389127808000} - \frac{147007268067241018958241836628521 n^{2}}{93256746094512000} + \frac{299407299371860260636969629 n}{44790183900} - 13783108115112102.
\]

\section{The count is a walk count}

The entry's defining condition is local: it is a condition on a bounded window of consecutive lines. Take as vertices the
admissible configurations of such a window and put an edge from $u$ to $v$ when $v$ may
follow $u$, which the condition decides by inspection. An admissible array is then exactly a
walk, so with $M$ the adjacency matrix of that digraph on $S=1992$ vertices, and $u,w$ the
vectors recording which windows may start and end an array,
\[
a(n)\;=\;u^{\!\top}M^{\,g(n)}w
\]
for the linear function $g$ that the entry's shape supplies. In particular $a$ is
$C$-finite: it satisfies some monic integer linear recurrence, though not necessarily a short
one.

\section{A closed form is a recurrence}

\begin{lemma}\label{lem:ann}
Let $f(m)=\sum_{j}p_j(m)\lambda_j^{\,m}$ with the $\lambda_j$ distinct and the $p_j$
polynomials, and put $q(x)=\prod_j(x-\lambda_j)^{\deg p_j+1}$. Then $q$ applied to the shift
operator annihilates $f$: writing $q(x)=x^{r}-\sum_{i=1}^{r}c_ix^{r-i}$,
\[
f(m)\;=\;\sum_{i=1}^{r}c_i\,f(m-i)\qquad\text{for every }m.
\]
\end{lemma}

\begin{proof}
The shift operator $E$ acts on $m\mapsto\lambda^{m}p(m)$ as
$(E-\lambda)\bigl(\lambda^{m}p(m)\bigr)=\lambda^{m+1}\bigl(p(m+1)-p(m)\bigr)$, and
$p(m+1)-p(m)$ has degree one less than $p$. So $(E-\lambda)^{\deg p+1}$ kills
$\lambda^{m}p(m)$, and the factors of $q$ commute.
\end{proof}

Here $q(x)=(x-1)^{31}$, of degree $31$; the closed form is a polynomial in $n$, so this is $(x-1)^{31}$, and the recurrence it encodes is
\begin{equation}\label{eq:rec}
a(n)\;=\;31\,a(n-1) - 465\,a(n-2) + 4495\,a(n-3) - 31465\,a(n-4) + 169911\,a(n-5) - 736281\,a(n-6) + 2629575\,a(n-7) - 7888725\,a(n-8) + 20160075\,a(n-9) - 44352165\,a(n-10) + 84672315\,a(n-11) - 141120525\,a(n-12) + 206253075\,a(n-13) - 265182525\,a(n-14) + 300540195\,a(n-15) - 300540195\,a(n-16) + 265182525\,a(n-17) - 206253075\,a(n-18) + 141120525\,a(n-19) - 84672315\,a(n-20) + 44352165\,a(n-21) - 20160075\,a(n-22) + 7888725\,a(n-23) - 2629575\,a(n-24) + 736281\,a(n-25) - 169911\,a(n-26) + 31465\,a(n-27) - 4495\,a(n-28) + 465\,a(n-29) - 31\,a(n-30) + a(n-31).
\end{equation}

\section{The proof}

\begin{lemma}\label{lem:crit}
With $q$ as above, $a$ satisfies \eqref{eq:rec} for all large $n$ if and only if
$u^{\!\top}M^{\,j}q(M)w=0$ for every $j\ge0$, and it is enough to check $j<S$.
\end{lemma}

\begin{proof}
Substituting the walk expression, the residual $a(n)-\sum_ic_ia(n-i)$ equals
$u^{\!\top}M^{\,g(n)-r}q(M)w$ up to the fixed linear reparametrisation $g$. For the bound,
$M^{S}$ is an integer combination of $I,M,\dots,M^{S-1}$ by Cayley--Hamilton, so the values
$u^{\!\top}M^{j}q(M)w$ satisfy the monic recurrence given by the characteristic polynomial of
$M$; $S$ consecutive zeros therefore force all later ones, and the last nonzero value pins the
threshold exactly.
\end{proof}

\begin{theorem}
$a(n)=f(n)$ for every $n\ge36$.
\end{theorem}

\begin{proof}
The vector $w'=q(M)w$ was formed by $31$ matrix--vector products in exact integer
arithmetic, and $u^{\!\top}M^{j}w'$ evaluated for $j=0,1,\dots$, again exactly. Those integers
vanish from the point corresponding to $n=67$ onwards, and the run was continued until the
working vector was identically zero, which settles every larger $j$ at once. So $a$ satisfies
\eqref{eq:rec} for every $n>66$. By Lemma~\ref{lem:ann} $f$ satisfies \eqref{eq:rec}
identically. Both were then evaluated in exact arithmetic at the $31$ consecutive indices
$n=67,\dots,97$ and agree at each. A monic recurrence of order $31$
determines $a(m)$ from $a(m-1),\dots,a(m-31)$, and for every $m\ge98$ both
$m>66$ and $m-31\ge67$ hold, so induction on $m$ gives $a(m)=f(m)$ for every
$m\ge67$.
Below $n=67$ the recurrence argument gives nothing, since \eqref{eq:rec} is only
established for $a$ from $n=67$ on. The remaining indices $n=36,\dots,66$ were
therefore checked one at a time, directly against the walk count in exact integer arithmetic,
and agree.
\end{proof}

The range is tight in the sense that was checked: at $n=35$ the closed form and the sequence differ, so no larger range is available.

\section{Verification}

Four checks, each independent of the algebra above.

First, the walk model reproduces every one of the $20$ terms the entry publishes,
exactly, in integer arithmetic. This is what ties the model to the entry: the model is built
by reading the entry's English, and a misreading gives a different digraph and different
counts.

Second, the closed form was evaluated in exact rational arithmetic --- not floating point ---
at every published term from $n=36$ on, and agrees with the entry's own DATA at each.

Third, the annihilation test of Lemma~\ref{lem:crit} was carried out exactly, and the model
and the closed form were then compared directly at sixty consecutive indices beyond
$n=36$, far past where the proof already settles the matter.

Fourth, the closed form was deliberately perturbed --- by a constant and by a term linear in
$n$ --- and each perturbed form was rejected by the same comparison, so the test is not one
that anything would pass.

\begin{thebibliography}{9}
\bibitem{oeis} The OEIS Foundation, \emph{The On-Line Encyclopedia of Integer
Sequences}, \url{https://oeis.org}, 2026. Sequence A240335.
\bibitem{stanley} R.~P.~Stanley, \emph{Enumerative Combinatorics}, Volume 1, 2nd ed.,
Cambridge University Press, 2011. (Transfer-matrix method, Section 4.7; $C$-finite sequences,
Section 4.1.)
\bibitem{kp} M.~Kauers and P.~Paule, \emph{The Concrete Tetrahedron}, Springer, 2011.
(Closed forms and linear recurrences with constant coefficients.)
\bibitem{horn} R.~A.~Horn and C.~R.~Johnson, \emph{Matrix Analysis}, 2nd ed., Cambridge
University Press, 2013.
\end{thebibliography}

\end{document}
